\documentclass[10pt,a4paper]{article}
\usepackage[utf8]{inputenc}
\usepackage[T1]{fontenc}
\usepackage[italian,provide=*]{babel}
\usepackage{geometry}
\usepackage{enumitem}
\usepackage{hyperref}
\usepackage{xurl}
\usepackage[expansion=false]{microtype}
\geometry{top=1.45cm,bottom=1.45cm,left=1.55cm,right=1.55cm}
\setlist[itemize]{noitemsep,topsep=2pt,leftmargin=*}
\setlength{\parindent}{0pt}
\setlength{\parskip}{3pt}
\setlength{\emergencystretch}{2em}
\hypersetup{colorlinks=true,urlcolor=blue,linkcolor=black}
\newcommand{\code}[1]{\texttt{#1}}

\title{\vspace{-1.2cm}\textbf{Archivio XML dei CIG ANAC}\\[2pt]\large Relazione sintetica sulla preparazione documentale}
\author{Massimiliano Ricchiuti\\Elaborazione di testi e gestione documentale}
\date{A.A. 2025/2026}

\begin{document}
\maketitle
\vspace{-1.1cm}

\section*{1. Obiettivo e dataset}
Il progetto realizza una pipeline Python che processa fonti relative a contratti pubblici, le trasforma in XML, valida i documenti mediante DTD, estrae informazioni e genera una collezione HTML navigabile. Il dataset comprende 15 XML, soglia prevista per il progetto individuale, derivati da fonti CSV, JSON, HTML e PDF conservate in \code{fonti\_originali/}. Il campione è didattico e non rappresenta l'intero universo ANAC.

La preparazione parte dai CIG presenti nei JSON, integra le righe omologhe del CSV e collega le ulteriori fonti. L'associazione non usa eccezioni dedicate: il programma cerca un CIG noto nel nome del file e, se il nome è descrittivo, nel testo estraibile. In questo modo il documento \emph{pn vsf15-25-sua lettera inv.-disciplinare indifferenziata 2025-26 Frascati} viene associato a \code{B6DD95EE23} perché il codice compare nel testo, insieme al PDF ANAC canonico. Lo stesso criterio si applica a nuovi documenti e ad altri CIG.

\section*{2. Modello XML e validazione}
Ogni radice \code{contratto} reca il CIG e contiene fonti, informazioni di gara, stazione appaltante, pubblicazioni, partecipanti, importi, date, CPV e descrizione documentale. Per ciascuna fonte sono conservati tipologia, MIME type, percorso e descrizione; le pagine di dettaglio propongono i relativi download.

La DTD \code{schema/contratto\_cig.dtd} stabilisce ordine, cardinalità e attributi. Il requisito del contenuto testuale è soddisfatto da un content model misto conforme alla sintassi XML:
\begin{verbatim}
<!ELEMENT descrizioneDocumentale
  (#PCDATA | tipologia | oggetto | categoria | termine | luogo)*>
\end{verbatim}
L'elemento alterna testo ed elementi semantici ripetibili; il dataset non è quindi composto soltanto da blocchi strutturati. La build valida tutti gli XML con \code{lxml}; un documento mal formato o non valido interrompe l'esecuzione e viene registrato in \code{output\_data/validation.xml}.

\section*{3. Estrazione, analisi e output}
Gli XPath convertono gli XML validi in istanze del modello intermedio \code{ContractRecord}. Su tale modello sono calcolate distribuzioni territoriali, copertura dei campi, sequenze cronologiche, statistiche economiche e frequenze testuali. I PDF sono esaminati con \code{pypdf}, con ripiego su \code{pdfminer.six}; gli importi usano \code{Decimal} e le incongruenze temporali restano visibili senza correzioni arbitrarie.

Jinja2 genera \code{index.html}, un archivio filtrabile, una pagina per ciascun XML valido, l'analisi del campione, la qualità dei dati e la pagina unificata \emph{Progetto e metodo}. Quest'ultima incorpora metodologia e documentazione, sostituendo le precedenti voci ambigue “Metodologia” e “PDF”. Ogni dettaglio contiene parte del contenuto elaborato, il download dell'XML e delle fonti collegate. Tutte le pagine di dettaglio recano microdata \code{schema.org/GovernmentService}; CSS, JavaScript e grafici SVG sono locali e responsivi.

\newpage
\section*{4. Riproducibilità, verifiche e uso di strumenti AI}
I comandi principali sono:
\begin{verbatim}
python -m pip install -r requirements.txt -r requirements-dev.txt
python scripts/build_all.py
python -m pytest
\end{verbatim}
La suite verifica parsing, DTD, estrazione, statistiche, mapping dei documenti descrittivi, pagine, microdata e collegamenti interni. GitHub Actions ripete build e test a ogni push su \code{main} e pubblica \code{dist/} su GitHub Pages.

Gli strumenti AI sono stati impiegati per la revisione architetturale, l'implementazione e il controllo della documentazione. Le proposte sono state sottoposte a validazione DTD, test automatici, build completa e ispezione dell'output. I prompt e le finalità d'uso sono riportati in \code{report/prompts\_utilizzati.md}. Le fonti originali, gli script di preparazione, la DTD, gli XML e gli output HTML restano disponibili e ispezionabili.

\section*{5. Riscontro dei requisiti}
\begin{itemize}
  \item 15 file XML distinti e validi, con directory di input configurabile;
  \item DTD creata per il progetto e content model misto effettivamente usato;
  \item fonti eterogenee CSV, JSON, HTML e PDF, con script di preparazione consegnato;
  \item una pagina HTML per ogni XML, indice navigabile e pagine aggiuntive di analisi;
  \item download dell'XML e dei documenti originali collegati al relativo CIG;
  \item foglio CSS comune, layout responsivo e annotazioni microdata generate;
  \item relazione entro tre pagine e prompt conservati nella documentazione.
\end{itemize}

\section*{6. Limiti}
Il dataset è circoscritto e presenta campi mancanti o semanticamente non omogenei; perciò le analisi sono descrittive. Il riconoscimento automatico richiede che il CIG sia presente nel nome o nel testo estraibile: per PDF costituiti esclusivamente da immagini è necessario un livello OCR. Tali limiti sono dichiarati nel sito e non vengono occultati mediante dati inventati.

\end{document}
